\chapter{Concept and Requirements}\label{ch:concept}

\section{Signal Path}
The input is stereo or mono and the output stereo, single-ended at every jack. The IN jack is TRS, tip left and ring right. The tip passes the relay into the input buffer and the codec's left ADC; the ring goes through its own buffer to the codec's right ADC. A mono plug grounds the ring, the right channel reads silence and the firmware copies left to right. The H750 runs the chord tracker and the string bank and mixes dry and wet. The codec's two DAC channels feed the left and right output buffers, and both outputs return through the relay. An expression pedal on a TRS jack goes through its own buffer to an MCU ADC input.

{\small
\begin{verbatim}
IN -> relay -> input buffer -> ADC (codec L) -> H750
     H750: front end -> chord tracker -> tuning manager
                     -> string bank x24 -> stereo spread
           dry + wet mix -> DAC L -> out buffer -> relay -> OUT L
                         -> DAC R -> out buffer -> relay -> OUT R
EXP (TRS) -> buffer -> H750 ADC
\end{verbatim}
}

The dry signal is converted too. The point of the pedal is the wet sound, and keeping the dry in the DSP puts the mix, the trails and the output level under firmware control. Both codec ADC channels are used. The AK4621EF has differential inputs only, so each channel has a THS4522 single-ended to differential driver, as on the Timekeeper.

\section{Engine}
\begin{table}[H]
  \small\renewcommand{\arraystretch}{1.12}
  \caption{Engine blocks and what each reuses.}\label{tab:engine}
  \begin{tabularx}{\textwidth}{@{}>{\raggedright\arraybackslash}p{28mm} L >{\raggedright\arraybackslash}p{26mm}@{}}
    \toprule
    \fmtophead{Block} & \fmtophead{What it does} & \fmtophead{Reuse} \\
    \midrule
    String bank & 24 Karplus--Strong / waveguide strings, each a fractional delay line with a one-pole loss filter and a tuning allpass. The guitar reaches each string through a narrow resonant band-pass at the string's pitch, so a string rings only when its note or an overtone is played. No envelope trigger. & New \\
    Jawari & Nonlinear bridge: a soft clip and a short extra delay tap in the loop, depth on the Jawari knob. Zero is a harp, full is a sitar. & New \\
    Chord tracker & FFT of the input, chroma vector and template match to a chord set with hysteresis. Only chords lasting more than about 400\,ms retune the strings. & Timekeeper spectral bank \\
    Tuning manager & Sets the string pitches from the chord (Follow), a fixed key (Key) or an open tuning (Drone, stacked fifths among the sets). Snap retunes only silent strings; Glide slides sounding strings to the new pitch; Lock stops retuning. & New \\
    Stereo spread & Strings panned by pitch, low to the left and high to the right. & New \\
    Front end & Hum filter and noise gate before the strings. & Timekeeper \\
    Framework & Effect chain, settings in QSPI flash, control smoothing, expression, UI simulator. & Timekeeper \\
    \bottomrule
  \end{tabularx}
\end{table}

Twenty-four strings at 44.1\,kHz cost a few MIPS each, and the Timekeeper's spectral bank alone is heavier, so the load should be comfortable on the H750 at 480\,MHz. It has \textbf{not been measured}. Measuring it on the Timekeeper board is the first step of the firmware plan (Chapter~\ref{ch:firmware}).

\section{Controls}
The face follows The Alchemist's grid (20\,mm columns at $-20$, 0 and $+20$, rows at $+38$ and $+13$, the same 20 by 25\,mm pitch as The Gremlin and the Timekeeper) with one more toggle at (0, $-5$), and 12.5\,mm knobs. Nothing through-hole sits over the upper jacks (Section~\ref{sec:toggleconflict}). The three toggles are Taiway 100-DP6 ON-ON-ON, The Alchemist's part, so each one stages the effect in three steps. Coordinates are face X and Y in millimetres from the enclosure centre.

\begin{table}[H]
  \small\renewcommand{\arraystretch}{1.12}
  \caption{Face controls.}\label{tab:controls}
  \begin{tabularx}{\textwidth}{@{}>{\raggedright\arraybackslash}p{30mm} >{\raggedright\arraybackslash}p{32mm} L@{}}
    \toprule
    \fmtophead{Position} & \fmtophead{Control} & \fmtophead{Function} \\
    \midrule
    ($-20$, $+38$) knob & Mix & Dry/wet \\
    (0, $+38$) knob & Sustain & String decay time \\
    ($+20$, $+38$) knob & Strings & 6 to 24 active strings \\
    ($-20$, $+13$) toggle & Tuning & Follow / Key / Drone \\
    (0, $+13$) knob & Jawari & Bridge buzz \\
    ($+20$, $+13$) toggle & Brightness & Dark / Warm / Glassy loss-filter voicings \\
    (0, $-5$) toggle & Retune & Snap / Glide / Lock \\
    band at $-18$ & 0.91\,in OLED strip & Note names the strings are tuned to, chord, mode and the value of the knob last moved \\
    ($-20$, $-35$) & Effect LED & \\
    ($-20$, $-49$) & Bypass footswitch & Momentary SPST-NO soft-touch on hand-solder pads \\
    ($+20$, $-49$) & Hold footswitch & Freezes the tuning and sets infinite sustain \\
    right wall & IN, EXP & Alchemist jack positions \\
    left wall & OUT L, OUT R & Alchemist jack positions \\
    top wall & DC & Alchemist DC position \\
    \bottomrule
  \end{tabularx}
\end{table}

All four pots and the three toggles go to the MCU; none is in the audio path. Without a screen, secondary settings (key, drone set, bypass mode True or Trails, expression target Bend or Swell, spread width and others) are reached by a footswitch gesture and set with the knobs, with the OLED strip as the display.

\section{Bypass Behaviour}
There is one relay, with true bypass on OUT~L. De-energised, the relay connects OUT~L to IN and grounds the input buffer, so the pedal passes signal unpowered; \textbf{OUT~R is silent in true bypass}. In Trails mode the relay stays energised and the DSP passes the dry signal to both outputs while the strings ring out, which is the stereo player's bypass. Before every relay change the DSP fades its outputs to silence, switches, waits for the contacts to settle and fades back in. There is no mute transistor.

\section{Requirements Baseline}
The firmware requirements are drafted in \texttt{docs/firmware-requirements.md} (Draft~A, 2026-10-04) with numbered, verifiable requirements. Table~\ref{tab:reqgroups} summarises the groups; Chapter~\ref{ch:firmware} describes the design they imply. Values marked TBC in that document are proposals to be settled on the Timekeeper prototype.

\begin{table}[H]
  \small\renewcommand{\arraystretch}{1.12}
  \caption{Firmware requirement groups.}\label{tab:reqgroups}
  \begin{tabularx}{\textwidth}{@{}>{\raggedright\arraybackslash}p{26mm} L@{}}
    \toprule
    \fmtophead{Group} & \fmtophead{Covers} \\
    \midrule
    PLT & Platform: H750 at 480\,MHz, internal RAM only, 128\,KB flash budget, one source tree for the Harp and Timekeeper boards, host-testable modules, no dynamic allocation \\
    AUD & Audio path: AK4621EF over SAI1 at 44.1\,kHz / 24-bit, stereo in (mono plug copied), dry at unity, latency target 3\,ms, wet limiter, no clicks \\
    FE, STR, JAW & Front end, string bank (24 strings, C2 to C6, $\pm$2 cents), Jawari \\
    CHD, TUN & Chord tracker (chord set, 400\,ms commit, 90\,\% accuracy target on clean DI) and tuning manager (three modes, Snap / Glide / Lock, Hold) \\
    OUT & Stereo spread, mix law, bypass-matched level \\
    POT, TGL, FSW, EXP, SEC & Controls, footswitches, expression and the secondary layer \\
    DSP & OLED strip and effect LED \\
    BYP & True and Trails bypass, the fade--switch--fade sequence, power-up state \\
    PER, SYS & Settings in QSPI flash, start-up, watchdog and fault handling \\
    PRF, ARC, VER & CPU, RAM and flash budgets, code structure and reuse, host tests and bench verification \\
    \bottomrule
  \end{tabularx}
\end{table}
